\documentclass[11pt]{article}

% Change "review" to "final" to generate the final (sometimes called camera-ready) version.
% Change to "preprint" to generate a non-anonymous version with page numbers.
\usepackage[review]{acl}

% Standard package includes
\usepackage{times}
\usepackage{latexsym}

% For proper rendering and hyphenation of words containing Latin characters (including in bib files)
\usepackage[T1]{fontenc}
% For Vietnamese characters
% \usepackage[T5]{fontenc}
% See https://www.latex-project.org/help/documentation/encguide.pdf for other character sets

% This assumes your files are encoded as UTF8
\usepackage[utf8]{inputenc}

% This is not strictly necessary, and may be commented out,
% but it will improve the layout of the manuscript,
% and will typically save some space.
\usepackage{microtype}

% This is also not strictly necessary, and may be commented out.
% However, it will improve the aesthetics of text in
% the typewriter font.
\usepackage{inconsolata}

%Including images in your LaTeX document requires adding
%additional package(s)
\usepackage{graphicx}

% If the title and author information does not fit in the area allocated, uncomment the following
%
%\setlength\titlebox{<dim>}
%
% and set <dim> to something 5cm or larger.

\title{What LLMs Can't See About Gender: A Close Reading of Evaluation Techniques}

% Author information can be set in various styles:
% For several authors from the same institution:
% \author{Author 1 \and ... \and Author n \\
%         Address line \\ ... \\ Address line}
% if the names do not fit well on one line use
%         Author 1 \\ {\bf Author 2} \\ ... \\ {\bf Author n} \\
% For authors from different institutions:
% \author{Author 1 \\ Address line \\  ... \\ Address line
%         \And  ... \And
%         Author n \\ Address line \\ ... \\ Address line}
% To start a separate ``row'' of authors use \AND, as in
% \author{Author 1 \\ Address line \\  ... \\ Address line
%         \AND
%         Author 2 \\ Address line \\ ... \\ Address line \And
%         Author 3 \\ Address line \\ ... \\ Address line}

\author{First Author \\
  Affiliation / Address line 1 \\
  Affiliation / Address line 2 \\
  Affiliation / Address line 3 \\
  \texttt{email@domain} \\\And
  Second Author \\
  Affiliation / Address line 1 \\
  Affiliation / Address line 2 \\
  Affiliation / Address line 3 \\
  \texttt{email@domain} \\}

%\author{
%  \textbf{First Author\textsuperscript{1}},
%  \textbf{Second Author\textsuperscript{1,2}},
%  \textbf{Third T. Author\textsuperscript{1}},
%  \textbf{Fourth Author\textsuperscript{1}},
%\\
%  \textbf{Fifth Author\textsuperscript{1,2}},
%  \textbf{Sixth Author\textsuperscript{1}},
%  \textbf{Seventh Author\textsuperscript{1}},
%  \textbf{Eighth Author \textsuperscript{1,2,3,4}},
%\\
%  \textbf{Ninth Author\textsuperscript{1}},
%  \textbf{Tenth Author\textsuperscript{1}},
%  \textbf{Eleventh E. Author\textsuperscript{1,2,3,4,5}},
%  \textbf{Twelfth Author\textsuperscript{1}},
%\\
%  \textbf{Thirteenth Author\textsuperscript{3}},
%  \textbf{Fourteenth F. Author\textsuperscript{2,4}},
%  \textbf{Fifteenth Author\textsuperscript{1}},
%  \textbf{Sixteenth Author\textsuperscript{1}},
%\\
%  \textbf{Seventeenth S. Author\textsuperscript{4,5}},
%  \textbf{Eighteenth Author\textsuperscript{3,4}},
%  \textbf{Nineteenth N. Author\textsuperscript{2,5}},
%  \textbf{Twentieth Author\textsuperscript{1}}
%\\
%\\
%  \textsuperscript{1}Affiliation 1,
%  \textsuperscript{2}Affiliation 2,
%  \textsuperscript{3}Affiliation 3,
%  \textsuperscript{4}Affiliation 4,
%  \textsuperscript{5}Affiliation 5
%\\
%  \small{
%    \textbf{Correspondence:} \href{mailto:email@domain}{email@domain}
%  }
%}

\begin{document}
\maketitle
\begin{abstract}
This paper applies theorizing from the Cultural Studies field of Trans Studies to analyze three recent Natural Language Processing (NLP) methods for detecting and evaluating gender bias in Large Language Models (LLMs). Drawing in particular on Trans Studies’ emphasis on gender as a visible, expressive, and social phenomenon, the paper examines how treating gender primarily as an internal identity limits what bias detection and evaluation methods are able to perceive. Methods that conceptualize gender as internal overlook external  elements of behavior and presentation only when they concern those groups that are most frequently targeted by discrimination. The paper concludes by outlining a direction for future research that integrate theoretical insights from Trans Studies to better capture the material and relational dimensions of gender bias in language technologies.

\end{abstract}

\section{Introduction}
There is a global urgency to define gender. In various countries around the world, gender designations like “transgender” and “nonbinary” are prompting backlash that has resulted in the retraction of hard-gained rights for LGBTQ+ groups, such as the recent ban on gender-affirming care for minors in the USA \citep{karni:2025}. This backlash takes what it terms “gender ideology”, or nontraditional views of gender that are not based on biological sex, as a rallying cry for political priorities that align with patriarchal and xenophobic logics, such as anti-immigrant discrimination (8, 25, 36).

In response to attacks on so-called gender ideology, some progressive groups have doubled down on a definition of gender that is based on internal identification rather than biological reality (16). Drawing from foundational theories of gender, such as Judith Butler’s (7) theory of Gender Performativity, which breaks down the rigid and causal relationship between sex and gender (in which a male sale denotes a male gender), this perspective leaves room for individuals to identify according to their personal perception and feeling, which may or may not correspond to their assigned sex at birth. However, as I explain below, this understanding of gender as internal identification mis-applies Butler’s critique, which explicitly situates gender as a behavioral phenomenon.

Building on the notion of gender as external, the academic field of Trans Studies explores paradigms that describe non-normative forms of gender according to its behavioral and expressive characteristics. One perspective from this field, called “trans materialism”, claims that gender ought to be defined by behavior, social relations, and appearance \citep{amin:2022}. Its reasoning states that if gender is conceived as an external phenomenon, it will be easier to make visible those who are non-conforming to gender norms, and therefore most marginalized.

Across the disciplinary divide, NLP has been developing various methods for studying gender bias in language forms (6, 35, 27). In the last few years, a growing number of these methods address queer and trans bias in particular (19, 33, 31, 15). This paper explores how theorizing from Trans Studies might be applied to three recent frameworks for detecting trans bias in language: TIDEs, EQUITBL, and TANGO (31, 15, 33). A close reading of each of these three frameworks examines in detail how gender is operationalized and how language makes gender forms visible. It argues that two of these frameworks, TIDEs and EQUITBL significantly benefit from addressing gender as a behavioral and social phenomenon, as theorized in Trans Studies. Conversely, the third approach, TANGO, is found to be limited by its focus on gender as primarily an internal quality. The discussion thus illustrates how Trans Studies’ theorizing on expressive and behavioral aspects of gender identity has direct application to techniques for marking up and coding gender in NLP. Answering calls for interdisciplinary theorizing in NLP (43, 4), this paper thus contributes to the increasing momentum around deconstructing “gender” as a concept within NLP by drawing from knowledge domains in Queer and Trans Studies (23, 9).

\section{Gender Bias in NLP}
Current techniques for detecting and mitigating gender bias often assume that gender equality means that gender terms ought to be interchangeable (i.e., switching “male” for “female” in text generation tasks). Word embedding (6) and Coreference Resolution tasks (35), for example, are based on the idea that each of the genders ought to have equal associations with all attributes (i.e. that a male gendered term ought to be as equally associated as a female one to attributes such as “babysitter” or “computer programmer”) (6). Meanwhile, other research (5, 21, 23, 9) points out that such frameworks do not consider how gender terms contain deeper associations with social norms. 

One major assumption that guides gender bias research is the gender binary itself, the assumption of gender as male or female (6, 22, 21). This tendency has been strongly critiqued in work that emphasizes LGBTIQIA+ identities and which, increasingly, centers these communities in the creation of bias detection and mitigation techniques (19, 14, 31, 33). From this research, one major trend is relevant to this project: the interrogation of how key terms like “gender” are being defined and operationalized. (23, 9), for example, draw from the insights of seminal gender theorist, Judith Butler, in order to conceptualize gender as a social construction. 

\section{Defining Gender}

The next subsections consider how gender has been conceptualized between popular discourses and the academic field of Trans Studies.

\subsection{Gender Ideology in Popular Discourse}
The term “gender ideology”, which the Southern Poverty Law Center classifies as a derogatory term, refers to the view that trans identities are internal and therefore false belief systems (37). In various countries around the world, gender ideology has risen to a national issue, serving to unite various political aims into what are mostly alt-right agendas. In the European continent, conservative groups have used gender ideology to prop up other forms of social discrimination, such as that against migrants. In Italy, for example, discourses of gender ideology have been linked to white nationalism, whereby “‘gender’ can be used to ‘stick’ ‘migrant’ and queer subjects together, characterizing both as threats to the sexual order of Europe” (26). Meanwhile, on the African continent, similar arguments are used toward politically opposite messaging, in anti-colonial discourses, where gender ideology is framed as a dangerous, Western import that threatens traditional gender roles. In Uganda, in particular, resistance against gender ideology has prompted a wave of anti-LGBT+ legislation, with one law penalizing what it calls “aggravated homosexuality” with the death penalty (2). Across these various political contexts, ranging from anti-immigrant to anti-West framings, gender ideology arguments tend to center on patriarchal arguments about the traditional sexual order.

In the USA, concerns over gender ideology have prompted a dramatic increase in legislation and policies which limits the rights of trans people (41, 39, 40, 38). Recent White House Executive Orders, for example, deploy a common, right-wing talking point that so-called “biological” women need explicit protection from trans women (40, 41). Other rules, for example, implicitly empowers federal officials like Immigration and Customs Enforcement (ICE) to detain persons that they believe to be trans, effectively expanding current practices of racial profiling to also include gender profiling (42). Turning on patriarchal and xenophobic logics similar to those in Europe and Africa, these initiatives illustrate that policies against gender ideology are in tandem with those reinforcing sex categories and immigration. 

These attacks on trans identities and rights are based on what the conservative side views as illegitimate conceptions of gender, particularly of gender as internal. For example, one of these Executive Orders describes gender ideology as “an ever-shifting concept of self-assessed gender identity, permitting the false claim that males can identify as and thus become women and vice versa, and requiring all institutions of society to regard this false claim as true” (39). This critique refers to the progressive view that gender identity is a subjective quality, an identification that is primarily internal rather than external. The Psychiatric Association, for example, defines gender identity as an “inner sense” of  “being a girl/woman, boy/man, some combination of both, or something else” (17). Similarly, the World Health Organization defines gender identity as “a person’s innate, deeply felt internal and individual experience of gender” (44).

However, some trans advocates argue that the progressive view weakens trans recognition and rights. Trans Studies scholars, in particular, contend that the concept of internalized gender inspires transphobic discrimination. For example, the progressive view of gender as an internal phenomenon influences transphobic reasoning on access to medical care, which questions why biological transition should be necessary if gender is an internal identification. Not only does this framing disempower those who seek medical transition, it also disregards other forms of transition, like name and pronoun changes, as part of trans identity. As others point out, many trans people choose not to undergo medical interventions and still identify as trans (29). The view of gender as internal, having to do with personal identity, thus opens the door to discrimination against multiple kinds of transition, including biological and social forms of transition.

According to this view, the right-wing push to disentangle gender from sex and to solidify sex as “binary and biological” is a reaction to the way that gender is conceptualized in trans affirming discourse on the left. One scholar, Andrea Long Chu \citet{chu:2024}, claims that “The left must reckon with its part in this. It has hung trans rights on the thin peg of gender identity”. Chu, a trans woman, explains that the progressive, pro-trans movement has “largely failed to form a coherent moral account of why someone’s gender identity should justify the actual biological interventions that make up gender-affirming care”.

\subsection{Gender Categories in Trans Studies}
The progressive view of gender descends from a widely accepted theory that gender is a social and behavioral phenomenon rather than a biological reality. This theory, developed by Judith Butler (7) as Gender Performativity theory in the early 1990s, dissolves the traditionally strict and causal relation between sex, gender, and sexuality. It stipulates that gender is not based on a biological reality, nor does it determine sexuality (e.g., with male gendered people preferring female gendered people). Rather, gender is an effect produced by social norms and expectations, which manifests as a series of behaviors (1991).

In the 35 years since Butler’s theory was introduced, the definition of gender has moved away from biology and into psychological concepts of identity. The primacy of behavior and social roles (which, Butler (8) is careful to point out, are not at will, but deeply compelled by societal rules and expectations), loses ground to internal concepts of identity and identification. Gender as identity, an internal feeling or identification as male, female, a combination of the two, or nonbinary genders, thus becomes separated from gender as expression, which denotes the masculine, feminine, or androgynous traits that are visibly apparent and enacted. Trans Studies scholar Kadji Amin (1) refers to this splitting of internal identity from external appearance or expression as the “divergence model”.  In the divergence model, when one term takes the normative position, another term absorbs the non-normative. According to Amin, the divergence model originally drives the split between the categories of “heterosexual” and “homosexual”, whereby “homosexual” was created first to encompass many variations of sex, gender, and sexuality non-conformity, while “heterosexual” was created by belatedly, in order to distinguish itself from “homosexual”. This view accords with that of sexual historiographers like Michel Foucault (20) and David Halperin (13) who, tracing the categories further back into the 19th century, explain that “heterosexual” as a category is a necessary contrast to that of “homosexual”, which was originally described as psychological pathology, such as in the work of early sexologists, Richard von Krafft-Ebing (30) and Havelock Ellis (17).

Subsequently, the divergence model which created the heterosexual category also created the cisgender category. As Amin (1) and Enke (18) argue, just as heterosexuality emerged as a response to homosexuality, so did the concept of cisgender, which responds to the emergence of the term “transgender” in the 1990s. Whereas the concept of homosexuality used to incorporate gender-variant forms of sexuality, for example, forms of male femininity such as the “queen”, and of female masculinity such as the “butch”, the term “transgender” offers a new category that can hold “all manner of gender-bending” (1). According to these scholars, it is not a coincidence that gay and lesbian rights begin to gain traction at the same time that the trans identity becomes solidified: as the trans category could now encompass non-normative gender expressions, it allows gender-conforming variants of “gay” and “lesbian” to attain more mainstream acceptance. 

The divergence model has, according to Amin (1), created a domino effect where the normative (cis)gender is increasingly idealized and, as a result, divorced from bodily expression. In other words, normative gender becomes both default and invisible—unmarked as such. Genders that comply with social expectations and norms will not be visible on the surface. Additionally, with each split, from homosexual/heterosexual to transgender/cisgender, the already internal notion of gender identity becomes buried “deeper and deeper within the private recesses of the self, where it increasingly disavows any relation to the social” (1). As an internal form, gender identity thus retains the privilege of being unnoticed from the outside.      

The corollary to this effect is that non-normative genders become hyper-visible. Because gender is regarded as internal, when dominant gender are invisible, those genders which are non-normative become marked and stigmatized. In other words, designating gender as an identity that is internal, that is not immediately apparent, effectively marginalizes any gender when it appears as an external expression. Amin (1) offers a common example from homosexual culture, where gay men who are normatively masculine, sometimes called “masc” or “butch” men, have higher social status than feminized gay men, “queens” and “fairies”, due to the cultural capital of masculinity. The more that non-normativity can be perceived in external ways, like the feminine appearance or mannerisms exhibited from a male-presenting person, the more that the person becomes recognized and stigmatized as other. Such stigmas are borne out in public life where, for example, masculine gay men are more likely to participate and achieve prominence over gay men who are visibly non-normative or gender-bending.

Thus, according to this view, gender as an internal identity reinforces normative expressions of gender externally. And as a result, visibly non-normative genders have a higher likelihood of experiencing discrimination for their visible difference. Given this, one way to resist the effects of gender stigmatization is to embrace gender as an external and relational phenomenon. In other words, defining gender as an external expression or behavior is more likely to make all genders visible. In the next section, I explore how three different gender bias frameworks might draw from the conception of gender as a visible and explicitly marked quality of identity.

\section{Methods for Detecting and Mitigating Gender Bias}
The paper now turns to the close reading of three bias detection and mitigation frameworks, which were chosen according to one central criteria: the incorporation of trans community input in the data collection, annotation, or evaluation stage. The close reading analysis emphasizes not only gender identities and expressions but also to language forms, such as tone and rhetorical context. As the paper explains in the conclusion, an analysis of gender in language needs to also consider language itself as an expressive form. This analysis is therefore guided by attention to the following: (1) gender as presentational, expressive, and/or behavioral; (2) elements of language like tone and rhetorical context and their relation to gender. 

\subsection{Framework 1: TIDEs}
The first bias detection framework is called Transphobia Identification in Digital Environments, or TIDEs, developed by Lameiro et al. (31). TIDEs comprises a dataset and a language model that detects trans bias in language. Building from the framework of a “trans technology,” a “technology that addresses the unique needs and challenges faced by trans people and communities” (24), this project draws from trans expertise for the data labelling work. The researchers aimed to create a dataset representing varied forms of transphobia, in particular, forms that are relatively more difficult to detect, like “TERFism”, i.e., “trans exclusionary radical feminism”, which identify as feminist despite negating the rights of trans women and girls. 

The data labelling takes into account not only the content of the sample, but also its tone and perspective. The careful data labelling work is aided by a detailed coding schema, which distinguishes between various kinds of transphobia, such as “misgendering”, “medical”, and “dogwhistles”, among others (31). This approach demonstrates that detecting transphobia requires attention not only to the content of the samples but also to cultural and rhetorical cues contained in language. For example, the following text is labelled as “not transphobic”, despite representing a transphobic point of view on the surface:
\begin{quote}
Hey guys, I just wanted to get your opinion on this:

Whenever I meet someone I really quickly try to establish whether or not they have a dick, I don’t just restrict this to potential sexual partners I mean everyone I meet. It just seems super important you know? Like how am I supposed to sleep at night unless I catalogue everyone’s genitalia?

(…) but here is the thing: people look at me like I’m crazy for asking! What the hell is up with that? I even had one girl tell me I was a totally creep!

Does this stuff happen to you?

/ok transjerkers I’m done for a little while, fuck it’s nice to vent
\end{quote}
Detecting this post correctly as “not transphobic” relies on recognizing its irony which, though suggested through its tone of enforced sincerity, is only confirmed at the end, with the assertion that “it’s nice to vent”. Additionally, the reference to “transjerkers”, an in-group moniker for a trans-dominated subreddit that emphasizes parody, also offers a clue to those who are already familiar with this online community. Therefore, the only hints that this text is not transphobic may be missed by those who lack the in-group cultural knowledge and language cues.

Contrast this example with one using a similar tone for a serious request. This post, which is also drawn from reddit, asks for advice on gender presentation:
\begin{quote}
Hey! I’m going through some gender experimentation, so I’d like to tap into masculinity for awhile(usually (sic) I identify as just pure androgynous, but I struggle with pulling that off). So, to begin with, I’ll try to master visual/behavioral androgyny first. I have long hair but don’t want to cut it(a (sic) lot of manly men have long hair anyway), so how could I fix it up to look more masculine? Also, how could I dress more masculine-like? I don’t have a very manly looking body shape, but I’d like to dress to give the illusion of one. Basically, any and all advice on how to appear androgynous/masculine is appreciated.
\end{quote}
The difference between this and the previous sample has to do with the rhetorical context. While the previous sample is written as parody, this one likely represents a serious request. The necessity of data labelers with insider knowledge thus becomes paramount.

Both of these samples also show how language related to physical aspects of gender presentation is prone to being read as discrimination. Despite their different valences, one ironic and the other genuine, both samples associate gender with physical presentation: in the first, there is a reference to genitalia, and in the second, references to gendered hair and dress. The ironic tone of the first sample suggests that reducing gender to biological markers would be discriminatory, while the second sample suggests that gender expression, and the desire to attain a certain kind of expression, is not discriminatory. 

\subsection{Framework 2: EQUITBL}
EQUITBL, the Explore, Query, and Understand Implicit Textual Bias in Language data project (15) uses a mixed-methods approach to analyze social bias in language, which deliberately addresses bias toward trans identities. 

Building from the authors’ previous work (14), EQUITBL expands its methodology to include text generation and close reading. In brief, the process involves prompting both English and Swedish language models using a set of templates designed to generate narratives about various identity groups. Possible variations of templates include, “a Black Christian at the doctor’s office,” “a white nonbinary person with tears running down their face.” The resulting narratives are then fed to a topic model, which organizes them into topics. Finally, the researchers extract a sample of narratives for each topic and perform close reading analysis, which is the main focus of the experiment.  

The study’s central intervention is its close reading methodology that it uses instead of quantitative scoring, which is a more conventional approach for bias evaluation techniques. As the researchers explain, they intentionally use qualitative metrics in order to “encourage users to critically engage with the topics and patterns… rather than merely being incentivized to chase ‘better’ scores” (15). This metric consists of a series of questions that the researcher uses to guide close reading, including “How are people and bodies described?”, as well as questions that interrogate gender’s relation to specific themes.

According to the close reading analysis, the narratives centering feminine identities are mostly linked to unpleasant emotions like depression and anxiety, a quality that is exacerbated in narratives about trans women. Narratives about masculine subjects, by contrast, are mostly positive in tone, associated with happy events like weddings and parties. Additionally, narratives about nonbinary persons express overt themes of struggle and the need for community support, which tend to end on a positive, “uplifting” note (15).

Amid these results, researchers note an unexpected pattern of refusal from the model, which declines to discuss certain topics and identity groups. In these refusals, the models claim that their programming instructions require them to either stay away from the topic of identity, which may promote “negative or derogatory language” (15). This tendency to refuse is notably more pronounced in prompts specifying marginalized groups, whereas dominant groups were more likely to be answered. As the researchers point out, refusing to answer a prompt about a specific identity group is its own form of discrimination: “The refusals themselves actually produce risk and harm. They suggest that the model likely cannot say anything nice, which is alarming when frequently repeated about minoritized groups” (15).

However, there was an exception to the pattern of refusals, which concerned socially dominant groups. When identity markers from socially dominant groups were combined in an additive way, such as “white”, “Christian”, and “man” to form “a white Christian man”, the model was more likely to refuse than similarly intersectional, though less dominant, identity combinations. For example, a model is more likely to refuse a prompt about a “white Christian man” than a “Black Christian man.” This quirk creates a strange effect where dominant groups were the most discussed, unless the identity of these groups were described in detail.

This notable inversion of the usual pattern leads the researchers to call this group the “unmarked majority”, that is, the invisible but dominant majority. A lack of identity specification in the model prompts will lead the model to assume a dominant identity which, researchers explain, “is assumed to be a man, as well as likely white, straight, cisgender, and Christian” (15). In an incisive speculation, the researchers posit that this “unmarked” quality of dominant identity groups explains why the model refuses to answer certain kinds of questions about them. Because such groups are “so often unmarked”, as the researchers explain, “specifying them draws extra attention to the concept of ‘identity,’ which the model has been discouraged from talking about” (15). 

\subsection{Framework 3: TANGO}
TANGO, (T)ow(A)rds centering tra(N)s(G)ender and nonbinary voices, is a dataset that centers TGNB, or trans and nonbinary, users in its goals and approach (33). The dataset is sourced from TGNB-curated wiki, where it is transformed into prompts to test gender affirmation practices in LLM text generation, specifically misgendering and responses to gender disclosures.

The methodology is mixed, involving both qualitative and quantitative steps. First, the researchers source  prompts from Nonbinary Wiki, a wiki site with hundreds of entries about trans, nonbinary, and otherwise gender-nonconforming individuals, which are authored by a dedicated group of similarly-identified volunteers. From the collection of entries, the researchers compile templates for two different types of prompts: misgendering and gender disclosure. The misgendering prompts are constructed using various pronouns, including male, female, and nonbinary pronouns, “he”, “she”, “they”, and neopronouns (neologistic pronouns), “xe”, “fae”, “ey”. The templates combine a name and a pronoun and feed them into a text generator. Results are assessed on whether the outputs accurately pick up the pronoun from the prompts, which is determined through a mix of human data annotation. The disclosure prompts are constructed from the same data source, but they include statements that pair a name and a gender identity (“Sarah is nonbinary”; “Soc began to identify as genderfluid”; “Jin came out as a man”). Results are assessed through Perspective API, an automated toxicity detector.

Results from the experiment find that notable levels of misgendering occur with the trans and nonbinary pronouns, “they”, in comparison to binary pronouns, “he”, “she”. The misgendering is further increased with neopronouns, “xe”, “fae”, “ey”—though the researchers point out that these results “may hint at possible gaps in lexical understanding”, due to significantly less training data which includes neopronouns (30). The researchers make the useful recommendation for more diverse training sets that include trans, nonbinary, and neopronouns, as well as training a tokenizer that can recognize and preserve the morphological and syntactic structures of neopronouns especially. 

Unlike the previous frameworks, this method conceptualizes gender in relation to the disclosure of gender identity. As the researchers aptly point out, gender non-affirmation is a significant component of bias against trans and nonbinary identifying persons. Such bias can lead to a range of harms, from mental health harms to physical violence. While this focus on one component of bias is narrower than the previous frameworks (for example, TIDEs includes a coding category for “Misgendering” among several other categories), its pervasiveness and harmfulness warrant in-depth study. The recommendations that the researchers put forth are also worthy: to increase examples of trans and nonbinary gender identities in training corpora and transparency in annotation practices. 
Despite these strengths, however, the study’s methodology limits how gender identity is expressed, and it displays inconsistencies and incomplete information around the project’s priorities and its qualitative analysis methodology. 

First, the study limits how gender identity is expressed in language forms; it frames gender as only related to a specific template for self-disclosure, “<subject> is <gender identity>”. The focus on gender according to this identification template perhaps reflects the priorities of the research team themselves, which mostly consist of Amazon Alexa developers, whose produce operates over voice control, where gender self-identification is more relevant than gender as a visual or behavioral phenomenon. 

Second, there is an inconsistency around the purported goal of centering TGNB voices. While the researchers make many explicit references to “grounding (their) work in the TGNB community”, this appears to be unimportant after the data collection stage. Here data annotators (hired from Amazon Mechanical Turk) are not questioned about their familiarity or positioning to TGNB identities (33). 

\section{Discussion}
The close reading of the TIDEs samples shows the importance of rhetorical context and tone for determining if language is transphobic or not. While the first example frames externalized elements of gender as grounds for discrimination, the second frames it as a means for expression and visibility. Both examples demonstrate \citet{amin:2022} theorizing on the importance of conceptualizing gender as external. Furthermore, they illustrate the deep understanding of community norms and rhetorical cues necessary for precise coding. 

The EQUITBL research also shows the importance of conceiving gender as an external phenomenon. In this study, attention to gender as external reveals implicit assumptions about which genders should be marked and which ones are not. The models are more likely to refuse prompts about “unmarked” groups (like “White” or “male” identities) because explicitly identifying these aspects is more conspicuous than identifying other groups. TANGO, meanwhile, considers gender in the context of identity self-affirmation or disclosure, which is comparably limited. As mentioned above, this focus on self-disclosure may be due to the commercial affiliations within the TANGO team which, as other research (4) has pointed out, may prioritize “performance” in narrow terms. 

Coming back to the TIDEs and EQUITBL studies, their focus on gender as a relational, contextual, and external phenomenon suggest that increasing the frequency of references to less dominant (marginalized) genders might be a viable technique for mitigating gender bias. Indeed, existing evaluations of LLMs show that the most frequent and default subject category is the straight white male, and that other subject categories, like woman, black and homosexual are less prevalent in both training data and model outputs (22). Additionally, research has also shown a correlation between the presence and frequency of certain identities in the training data and their association with positive and negative terms: van Loon et al. (32), in particular, find that the less frequently a group is named, the more negative the associations with that group. If a group’s relatively frequent presence in the training data has a positive effect on their representation, one might reasonably assume that making marginalized subjects more frequent in the training data, to be on par with the frequency of representation of dominant subjects, may mitigate bias about these groups in model outputs. 

However, any solution must also take into account that the default subject is often unmarked. As et al. (15) point out, the majority of the texts about white male subjects were the ones least likely to describe that subject in external, physical terms. That is perhaps because of an inverse relationship between physical markers of identity and positions of social dominance, which Devinney et al.’s (15) work bear out. Here, fewer physical descriptions themselves mark social dominance. In those results, white women are described in physical terms only slightly more than white men, while marginalized identities are described in more detail than both groups (15). These results suggest that the more that a person is described in physical terms, the more likely they come from a marginalized group. Increasing physical descriptions of marginalized subjects, then, may not help to mitigate negative associations with those groups. 

Given this complex dynamic between group visibility and dominance, this paper makes one main recommendation for future work on gender bias in NLP: that researchers increase physical descriptions of all gendered subjects in training sets, including those from dominant genders, especially white, male, heterosexual, and Christian identities. Some work already exists toward this goal. In bias mitigation research by Reif and Schwartz (34), “Fighting Bias with Bias,” the researchers share a promising approach: amplifying rather than reducing bias in a model’s training dataset. As they point out, “filtering (out bias) can obscure the true capabilities of models to overcome biases, which might never be removed in full from the dataset” (34). Instead of an approach that reduces bias, they offer one of increasing it, such as by including phrases like “the pretty doctor” so that a model will interpret “doctor” to be female. Future work might develop this practice to mark up targets and attributes which describe gender as physical and behavioral, rather than as internal or psychological, for all identity groups regardless of dominance.

\section{Conclusion}
This paper has shown that contemporary concepts of gender that frame it as an internal identity, detached from outward expression and embodiment, produce unintended consequences: idealized and invisible norms (such as cisgender and heterosexuality) are reinforced, while visibly non-normative gender expressions become further stigmatized. The problem with the universal, unmarked subject is precisely that, despite not being marked, he is overrepresented. This overrepresentation does not result from persistent or explicit labels, because the dominant subject does not carry a descriptive label as such: he simply is.


\section*{Limitations}

While this paper explores theoretical questions in existing NLP frameworks, it does not represent the full diversity of trans and gender-diverse experiences. The analysis is limited to interpretive readings of three evaluation frameworks, and any attempt to operationalize or standardize gender in computational systems risks oversimplification.

\section*{Acknowledgments}

% Bibliography entries for the entire Anthology, followed by custom entries
%\bibliography{custom,anthology-overleaf-1,anthology-overleaf-2}

% Custom bibliography entries only
\bibliography{custom}

\end{document}
